\documentclass[12pt,a4paper]{article}

% Biên dịch bằng XeLaTeX
\usepackage[a4paper,margin=2cm]{geometry}
\usepackage{fontspec}
\usepackage{polyglossia}
\setdefaultlanguage{vietnamese}
\setmainfont{Times New Roman} % Có thể đổi sang TeX Gyre Termes nếu máy không có Times New Roman

\usepackage{graphicx}
\usepackage{tikz}
\usepackage{array}
\usepackage{setspace}
\usepackage{calc}
\pagestyle{empty}
\setlength{\parindent}{0pt}

% ====== TÙY CHỈNH ĐƯỜNG DẪN LOGO TẠI ĐÂY ======
% Ví dụ: images/logo.png hoặc logo.png
\newcommand{\LogoPath}{images/logo.png}

% Khung viền xanh kiểu góc bậc như ảnh mẫu
\newcommand{\CoverFrame}{%
  \begin{tikzpicture}[remember picture,overlay]
    \draw[line width=1.2pt,color=blue!55!black]
      ([xshift=1.15cm,yshift=-1.15cm]current page.north west)
      rectangle
      ([xshift=-1.15cm,yshift=1.15cm]current page.south east);

    % Các nét vuông trang trí ở bốn góc
    \foreach \sx/\sy in {1/1,-1/1,1/-1,-1/-1}{
      \begin{scope}[shift={({\sx*1.15cm},{\sy*1.15cm})}]
      \end{scope}
    }
    \draw[line width=2.2pt,color=blue!55!black]
      ([xshift=1.15cm,yshift=-1.15cm]current page.north west) --
      ([xshift=1.42cm,yshift=-1.15cm]current page.north west) --
      ([xshift=1.42cm,yshift=-1.42cm]current page.north west);
    \draw[line width=2.2pt,color=blue!55!black]
      ([xshift=-1.15cm,yshift=-1.15cm]current page.north east) --
      ([xshift=-1.42cm,yshift=-1.15cm]current page.north east) --
      ([xshift=-1.42cm,yshift=-1.42cm]current page.north east);
    \draw[line width=2.2pt,color=blue!55!black]
      ([xshift=1.15cm,yshift=1.15cm]current page.south west) --
      ([xshift=1.42cm,yshift=1.15cm]current page.south west) --
      ([xshift=1.42cm,yshift=1.42cm]current page.south west);
    \draw[line width=2.2pt,color=blue!55!black]
      ([xshift=-1.15cm,yshift=1.15cm]current page.south east) --
      ([xshift=-1.42cm,yshift=1.15cm]current page.south east) --
      ([xshift=-1.42cm,yshift=1.42cm]current page.south east);

    % Các ô vuông nhỏ tạo hiệu ứng bậc ở góc
    \foreach \x/\y in {1.15/-1.15,-1.15/-1.15,1.15/1.15,-1.15/1.15}{
      \fill[blue!55!black]
        ([xshift=\x cm,yshift=\y cm]current page.north west) rectangle ++(0.10cm,-0.10cm);
    }
  \end{tikzpicture}%
}

\begin{document}
\CoverFrame

\begin{center}
  \vspace*{0.15cm}

  {\bfseries\large HỌC VIỆN CÔNG NGHỆ BƯU CHÍNH VIỄN THÔNG\par}
  \vspace{0.35cm}
  {\bfseries\large KHOA CÔNG NGHỆ THÔNG TIN 1\par}

  \vspace{0.55cm}
  {\large \rule{1.0cm}{0.4pt}\ \raisebox{-0.15em}{\(\blacksquare\)}\ \rule{1.0cm}{0.4pt}\par}

  \vspace{0.55cm}
  % Logo được lấy từ đường dẫn trong \LogoPath ở đầu file
  \IfFileExists{\LogoPath}{
    \includegraphics[width=3.5cm,height=3.5cm,keepaspectratio]{\LogoPath}
  }{
    \fbox{\parbox[c][3.2cm][c]{3.2cm}{\centering
      \small Chưa tìm thấy logo\\[2pt]
      \texttt{\LogoPath}}}
  }

  \vspace{0.35cm}
  {\bfseries\Large ĐỒ ÁN TỐT NGHIỆP\par}

  \vspace{0.45cm}
  {\bfseries\large\underline{Đề tài:}\par}

  \vspace{0.35cm}
  \begin{minipage}{0.88\textwidth}
    \centering
    \bfseries
    \large
    Nghiên cứu phát hiện đường chín đoạn trái phép trong\\
    hình ảnh bằng Faster R-CNN thông qua cân chỉnh biểu\\
    diễn thị giác và chắt lọc tri thức từ DINOv2
  \end{minipage}
\end{center}

\vspace{1.25cm}

\begin{center}
\begin{tabular}{@{}p{3.7cm}p{3.2cm}p{2.6cm}@{}}
  Giảng viên hướng dẫn: & TS. Đặng Hoàng Long & \\
  Nhóm sinh viên thực hiện: & Họ và tên & Mã sinh viên \\[0.12cm]
  & Đỗ Trần Tú & B22DCKH109 \\
  & Mai Danh Hiếu & B22DCKH039
\end{tabular}
\end{center}

\vfill

\begin{center}
  {\bfseries Hà Nội - 2026\par}
  \vspace*{0.55cm}
\end{center}

\end{document}
